\documentclass[11pt]{article}

\usepackage[margin=1in]{geometry}
\usepackage[T1]{fontenc}
\usepackage[utf8]{inputenc}
\usepackage{lmodern}
\usepackage{microtype}
\usepackage{amsmath,amssymb,amsthm,mathtools,bm}
\usepackage{booktabs,array,longtable,tabularx,multirow}
\usepackage{enumitem}
\usepackage{xcolor}
\usepackage{listings}
\usepackage{fancyhdr}
\usepackage{hyperref}
\usepackage[nameinlink,noabbrev]{cleveref}

\definecolor{deepblue}{RGB}{25,65,125}
\definecolor{deepgreen}{RGB}{25,105,70}
\definecolor{deepred}{RGB}{145,45,40}
\definecolor{softgray}{RGB}{246,247,249}
\definecolor{softblue}{RGB}{239,245,252}
\definecolor{softgreen}{RGB}{239,248,243}
\definecolor{softred}{RGB}{252,242,241}

\hypersetup{
  colorlinks=true,
  linkcolor=deepblue,
  citecolor=deepgreen,
  urlcolor=deepblue,
  pdftitle={Local Support, Kummer Coherence, and the Additive Skeleton},
  pdfauthor={OpenAI GPT-5.6 Pro}
}

\pagestyle{fancy}
\fancyhf{}
\fancyhead[L]{Alaoglu--Erd\H{o}s: local support and Kummer coherence}
\fancyhead[R]{22 August 2026}
\fancyfoot[C]{\thepage}
\setlength{\headheight}{14pt}

\newtheorem{theorem}{Theorem}[section]
\newtheorem{proposition}[theorem]{Proposition}
\newtheorem{lemma}[theorem]{Lemma}
\newtheorem{corollary}[theorem]{Corollary}
\newtheorem{conjecture}[theorem]{Conjecture}
\theoremstyle{definition}
\newtheorem{definition}[theorem]{Definition}
\newtheorem{problem}[theorem]{Research target}
\newtheorem{example}[theorem]{Example}
\theoremstyle{remark}
\newtheorem{remark}[theorem]{Remark}
\newtheorem{warning}[theorem]{Audit warning}

\newcommand{\Q}{\mathbb{Q}}
\newcommand{\Z}{\mathbb{Z}}
\newcommand{\R}{\mathbb{R}}
\newcommand{\F}{\mathbb{F}}
\newcommand{\ord}{\operatorname{ord}}
\newcommand{\Gal}{\operatorname{Gal}}
\newcommand{\rank}{\operatorname{rank}}
\newcommand{\vqp}{v_q}
\newcommand{\Span}{\operatorname{span}}
\newcommand{\dens}{\operatorname{dens}}
\newcommand{\supp}{\operatorname{supp}}
\newcommand{\lcm}{\operatorname{lcm}}
\newcommand{\AEconj}{\textsc{AE}}
\newcommand{\newresult}{\textcolor{deepgreen}{\textbf{NEW--UNCONDITIONAL}}}
\newcommand{\classical}{\textcolor{deepblue}{\textbf{CLASSICAL INPUT}}}
\newcommand{\computational}{\textcolor{deepred}{\textbf{COMPUTATIONAL CHECK}}}
\newcommand{\target}{\textcolor{deepred}{\textbf{OPEN TARGET}}}

\newcommand{\boxedparagraph}[2]{%
\begin{center}
\fcolorbox{#1}{softgray}{\begin{minipage}{0.92\textwidth}\small #2\end{minipage}}
\end{center}}

\lstdefinestyle{pythonstyle}{
  language=Python,
  basicstyle=\ttfamily\scriptsize,
  keywordstyle=\color{deepblue}\bfseries,
  commentstyle=\color{deepgreen},
  stringstyle=\color{deepred},
  numbers=left,
  numberstyle=\tiny\color{gray},
  stepnumber=1,
  numbersep=7pt,
  showstringspaces=false,
  breaklines=true,
  frame=single,
  backgroundcolor=\color{softgray},
  columns=fullflexible
}

\title{\textbf{Local Support, Kummer Coherence, and the Additive Skeleton}\\[4pt]
\large A Nonduplicative Progress Report on the Alaoglu--Erd\H{o}s Integer-Exponent Conjecture}
\author{OpenAI GPT-5.6 Pro}
\date{22 August 2026}

\begin{document}
\maketitle

\begin{abstract}
Let $x\in\R$ and suppose
\[
        M:=2^x\in\Z,\qquad A:=3^x\in\Z.
\]
The Alaoglu--Erd\H{o}s conjecture asserts that $x\in\Z$.  The supplied unified report already
contains an exceptionally broad investigation of transcendence, interpolation determinants,
$p$-adic amplification, finite differences, matrix and operator methods, Sturmian coding,
Pad\'e and canonical-product constructions, and many other routes.  This continuation therefore
does not repackage those methods.  It develops a different local--global program built around
the classical support problem for the multiplicative group and exact Kummer--Chebotarev
statistics.

The principal new result is an exact density theorem.  Put
$\Gamma=\langle 2,3,M,A\rangle\subset\Q_{>0}^{\times}$.  A nonintegral candidate has
$\rank\Gamma\in\{3,4\}$.  For an admissible odd prime $q$, an integer $e\ge1$, and
$Q=q^e$, a prime $p\equiv1\pmod Q$ supplies four $Q$-power residue logarithms.  We call
$p$ \emph{$Q$-coherent} when the output pair $(M,A)$ is a common scalar multiple of the
base pair $(2,3)$ in this residue quotient.  If the rank is four, the conditional density of
$Q$-coherent primes is
\[
 \frac{1+\dfrac{q(q^2-1)(q^{3e}-1)}{q^3-1}}{q^{4e}};
\]
if the rank is three and $q$ avoids one explicit determinant, it is
\[
 \frac{1+\dfrac{q(q^{2e}-1)}{q+1}}{q^{3e}}.
\]
Both are asymptotic to a constant times $q^{-e}$.  By contrast, an integral exponent is
$Q$-coherent at every unramified prime.  Thus a hypothetical nonintegral solution is
archimedean-coherent but, at arbitrarily deep prime-power residue levels, almost always
finite-place incoherent.  This is a sharp quantitative local signature, although it is not by
itself a contradiction.

Two complementary results are proved.  First, Kummer theory gives positive-density failures of
order support between $2$ and $M$ (and between $3$ and $A$) whenever the latter is not an
integral power of the former.  Coupled with the theorem of Corrales--Rodrig\'a\~nez and Schoof,
this identifies an exact sufficient target: prove that the candidate-induced monoid map on the
$3$-smooth numbers preserves congruences.  Second, every primitive binary relation
$u+v=w$ in that monoid is classified, and all same-time polynomial additive defects are shown
to generate only the single ideal
\[
       (M-2,A-3)=\gcd(M-2,A-3)\Z.
\]
The corresponding dynamic anchors
$\gcd(M^n-2^n,A^n-3^n)$ have subexponential size by the $S$-unit gcd theorem.  This gives a
rigorous ceiling for this strategy: same-time polynomial additive defects alone cannot
manufacture the exponential support needed by the local criterion.

The conjecture is not proved here.  The report instead isolates a new, exact density gap; gives
self-contained finite-ring counts and reproducible computations; identifies a finite regulator
or reciprocity statement that would bridge the remaining archimedean/local divide; and separates
Lean-friendly algebra from the external Chebotarev, support-theorem, and Subspace-Theorem inputs.
\end{abstract}

\tableofcontents
\newpage

\section{Scope, nonduplication, and claim ledger}

\subsection{What is inherited and what is new}

The attached source was decompressed and audited before this continuation was written.  It has
$20{,}925$ source lines and roughly $114{,}000$ words.  Its core and twenty-seven continuations
already treat, among other subjects:
\begin{itemize}[leftmargin=2em,itemsep=2pt]
\item equivalent four-exponentials formulations and the exact transcendence barrier;
\item rational-function rigidity, perfect powers, valuation lattices, and algebraic auxiliary
      bases;
\item second iterates and exponential towers;
\item arbitrary-node and semigroup Vandermonde determinants, min-plus valuations, transport
      ceilings, and finite determinant certificates;
\item Sturmian synchronization, Baker-type separation, sparse interpolation, random phases,
      moments, M\"untz approximation, Pad\'e defects, canonical products, and Loewner matrices;
\item structural-prime residuals, carry languages, operator coherence, and extensive Lean
      blueprints.
\end{itemize}
Those subjects are not redeveloped below.  Only three inherited facts are used as a baseline:
(1) a rational candidate exponent is integral; (2) the special four-exponentials instance is the
central archimedean wall; and (3) a nonintegral candidate has multiplicative rank three or four.
For completeness, the short proof of (3) is repeated in \cref{sec:rank}, because it is needed to
state the new density theorem without importing a hundred-page dependency chain.

The following notions and calculations form the new track of this report:
\begin{enumerate}[leftmargin=2.2em,itemsep=2pt]
\item the support preorder $u\preccurlyeq v$ and its exact application to $2,M$ and $3,A$;
\item $q^e$-power residue logarithms and a common-exponent coherence condition;
\item exact Chebotarev densities in both multiplicative-rank branches;
\item a finite residue-pattern detector for the unique rank-three relation;
\item the complete additive skeleton of the $3$-smooth monoid and its defect ideal;
\item the dynamic anchor sequence $C_n=\gcd(M^n-2^n,A^n-3^n)$;
\item a congruence-preservation target connected to Bhargava $p$-orderings;
\item a finite-regulator formulation of the missing archimedean/local bridge.
\end{enumerate}
A literal search of the supplied source found no occurrence of ``Chebotarev'',
``Corrales'', ``Schoof'', ``support problem'', ``discrete log'', or
``congruence-preserving''.  Kummer theory and the Bugeaud--Corvaja--Zannier gcd theorem do appear
there, but for different radical and cancellation diagnostics; the exact residue-density results
below do not.

\subsection{Status vocabulary}

Every substantial assertion is tagged conceptually as follows.
\begin{description}[style=nextline,leftmargin=1.8em]
\item[\newresult] A theorem proved in this report from stated inputs.  The finite-ring counts are
fully elementary; the density conversion uses standard Kummer theory and Chebotarev.
\item[\classical] A published theorem used as a black box: Gelfond--Schneider,
Chebotarev, Corrales--Rodrig\'a\~nez--Schoof, Bhargava's $p$-ordering framework, or the
Corvaja--Zannier $S$-unit gcd estimate.
\item[\computational] Exact modular arithmetic over a finite set of primes.  These checks test
formulas but are not evidence for the Alaoglu--Erd\H{o}s conclusion itself.
\item[\target] A statement whose proof would advance or settle the conjecture.  No target is
silently promoted to a theorem.
\end{description}

\boxedparagraph{deepred}{\textbf{Bottom line.}  No unconditional proof of the Alaoglu--Erd\H{o}s
conjecture is claimed.  The new achievement is an exact local density dichotomy and a precise
support-transfer target, together with a proof that the most obvious additive route to that
target collapses to a one-modulus invariant.}

\subsection{Audit of the reported 2026 breakthroughs}

The rival-team claims were checked only to calibrate the competitive context; none is used in a
mathematical argument below.
\begin{itemize}[leftmargin=2em]
\item An explicit polynomial map $\mathbb{C}^3\to\mathbb{C}^3$ with constant Jacobian determinant and a
three-point fibre was publicly circulated in July 2026 and independently formalized in Lean.
It refutes the Jacobian conjecture in dimensions at least three; the plane case remains open.
The relevant public artifacts are cited in the bibliography.
\item The ICARM elliptic-curve leaderboard records a curve submitted on 20 August 2026 with a
rank lower bound at least $30$ and a witness consisting of thirty independent rational points.
The leaderboard commentary attributes the construction to Claude with Levent Alp\"oge and Ava
Howell.  This is a certified lower bound, not an unconditional proof that the rank is exactly
$30$.
\item OEIS was updated in August 2026 to report Hadamard matrices for the twelve previously
unresolved admissible orders at most $2000$, attributing the constructions to Levent Alp\"oge,
Philippe Voinov, Saul Reynolds-Haertle, and Claude.  The existence checks themselves are finite
exact matrix computations.
\item No public preprint, repository, theorem statement, Lean file, or independently checkable
artifact for the rumored Alaoglu--Erd\H{o}s breakthrough was located as of 22 August 2026.
Accordingly, the rumor supplies neither a lemma nor a hint used here.
\end{itemize}

\section{Normalization and multiplicative rank}
\label{sec:rank}

Throughout,
\[
        2^x=M\in\Z_{>0},\qquad 3^x=A\in\Z_{>0}.
\]
If $x<0$, then $0<2^x<1$, impossible for a positive integer; if $0<x<1$, then
$1<2^x<2$, also impossible.  Thus either $x=0$ or $x\ge1$.

\begin{lemma}[Rational candidates]\label{lem:rational}
If $x\in\Q$ and $2^x\in\Z$, then $x\in\Z$.
\end{lemma}
\begin{proof}
Write $x=a/b$ with $b>0$ and $\gcd(a,b)=1$.  From $2^a=M^b$, comparison of the
$2$-adic valuations gives $a=bv_2(M)$, hence $b\mid a$ and $b=1$.
\end{proof}

\begin{corollary}[Transcendence of a counterexample]\label{cor:trans}
A nonintegral candidate $x$ is transcendental.
\end{corollary}
\begin{proof}
By \cref{lem:rational}, $x$ is irrational.  If it were algebraic, the
Gelfond--Schneider theorem would make $2^x$ transcendental, contrary to $2^x=M\in\Z$.
\end{proof}

Let
\[
   \Gamma=\langle 2,3,M,A\rangle\subset\Q_{>0}^{\times}
\]
and write $r=\rank\Gamma$.

\begin{proposition}[Rank dichotomy]\label{prop:rankdichotomy}
For a nonintegral candidate, $r\in\{3,4\}$.
\end{proposition}
\begin{proof}
Clearly $r\ge2$ because $2$ and $3$ are multiplicatively independent.  Suppose $r=2$.
Then $M$ and $A$ lie in the group generated by $2$ and $3$, so for integers
$a,b,c,d$,
\[
       M=2^a3^b,\qquad A=2^c3^d.
\]
Set $\theta=\log_2 3$.  The number $\theta$ is transcendental: it is not rational,
and if it were algebraic irrational then $2^\theta=3$ would contradict
Gelfond--Schneider.  Taking logarithms gives
\[
       x=a+b\theta=d+\frac{c}{\theta}.
\]
Hence
\[
       b\theta^2+(a-d)\theta-c=0.
\]
Transcendence of $\theta$ forces $b=c=0$ and $a=d$, so $x=a\in\Z$, a contradiction.
Thus $r\ne2$, while $r\le4$ by definition.
\end{proof}

If $r=3$, the relation lattice has rank one.  Fix its primitive generator
\[
  \bm c=(c_1,c_2,c_3,c_4)\in\Z^4,
  \qquad
  2^{c_1}3^{c_2}M^{c_3}A^{c_4}=1,
\]
unique up to sign, and define
\[
       \Delta(\bm c)=c_1c_4-c_2c_3.
\]

\begin{lemma}[Nondegenerate rank-three relation]\label{lem:delta}
For a nonintegral candidate in the rank-three branch,
$\Delta(\bm c)\ne0$.
\end{lemma}
\begin{proof}
If $\Delta=0$, the rational vectors $(c_1,c_2)$ and $(c_3,c_4)$ are proportional.
Neither pair can be zero, because $2,3$ are independent and because $M,A$ cannot satisfy a
nontrivial relation without forcing $x$ rational.  Thus there is $t\in\Q$ with
$(c_1,c_2)=t(c_3,c_4)$.  Taking logarithms of the multiplicative relation and using
$\log M=x\log2$, $\log A=x\log3$ gives
\[
  (t+x)(c_3\log2+c_4\log3)=0.
\]
The second factor is nonzero, so $x=-t\in\Q$, and \cref{lem:rational} yields
$x\in\Z$, contradiction.
\end{proof}

\subsection{Admissible residue primes}

The exact density formulas require exclusion of only finitely many primes.  This can be made
completely explicit from a valuation matrix.

Let $S$ be the finite set of rational primes dividing $6MA$.  Form the integer matrix $V$ whose
columns are the valuation vectors of $2,3,M,A$ at all primes in $S$; zero rows may be omitted.
Its rank is $r$.  Let $d_1\mid\cdots\mid d_r$ be the nonzero invariant factors in the Smith
normal form of $V$, and put
\[
       \sigma_\Gamma=d_1\cdots d_r.
\]
The number $\sigma_\Gamma$ measures the failure of the valuation lattice to be saturated.

\begin{definition}[Admissible prime]\label{def:admissible}
An odd prime $q$ is \emph{admissible for $\Gamma$} if
\[
      q\nmid 6MA\sigma_\Gamma,
\]
and, in the rank-three branch, also $q\nmid\Delta(\bm c)$.
\end{definition}
There are infinitely many admissible primes.

For $Q=q^e$, let $K_Q=\Q(\zeta_Q)$ and
\[
 L_Q=K_Q\bigl(2^{1/Q},3^{1/Q},M^{1/Q},A^{1/Q}\bigr).
\]
The next elementary saturation lemma is the algebraic engine behind all later densities.

\begin{lemma}[Exact Kummer relation module]\label{lem:kummermodule}
Let $q$ be admissible and $Q=q^e$.  For
$\bm n=(n_1,n_2,n_3,n_4)\in\Z^4$, the rational number
$2^{n_1}3^{n_2}M^{n_3}A^{n_4}$ is a $Q$th power in $K_Q$ if and only if
\[
 \bm n\in \mathcal R+Q\Z^4,
\]
where $\mathcal R$ is the integral relation lattice of $(2,3,M,A)$.  Consequently,
\[
 \Gal(L_Q/K_Q)\cong
 \begin{cases}
   (\Z/Q\Z)^4,&r=4,\\
   \{\bm t\in(\Z/Q\Z)^4:\bm c\cdot\bm t=0\},&r=3,
 \end{cases}
\]
and the latter group has order $Q^3$.
\end{lemma}
\begin{proof}
The reverse implication is immediate.  For the forward implication, assume the indicated
rational number is a $Q$th power in $K_Q$.  Because $q\nmid6MA$, every rational prime appearing
in the four generators is unramified in $K_Q/\Q$.  At any prime ideal above such a rational
prime $\ell$, the valuation of a $Q$th power is divisible by $Q$.  Therefore
\[
        V\bm n\in Q\Z^S.
\]
Since $q\nmid\sigma_\Gamma$, the Smith normal form shows that the valuation lattice is
$q$-saturated at every depth: the last congruence implies
$\bm n\in\ker(V)+Q\Z^4=\mathcal R+Q\Z^4$.

Because $K_Q$ contains all $Q$th roots of unity, ordinary Kummer theory identifies the Galois
group with the dual of the subgroup generated by the four classes in
$K_Q^\times/(K_Q^\times)^Q$.  The relation-module computation gives the displayed groups.  In
rank three, admissibility makes the primitive relation remain primitive modulo $q$, so its
kernel has cardinality $Q^3$.
\end{proof}

\begin{remark}
The use of explicit Kummer-degree theory can remove or optimize the finite exceptional set, but
it is not needed for the qualitative or exact formulas here.  The valuation-matrix definition is
algorithmic, elementary, and adequate for formalization.
\end{remark}

\section{The support problem as an exact bridge}
\label{sec:support}

The archimedean equations $M=2^x$ and $A=3^x$ do not themselves have a reduction modulo a
prime when $x$ is nonintegral.  The support problem supplies a different local object: the
orders of the rational integers $2,M$ and $3,A$ in finite fields.

\begin{definition}[Support preorder]
For $u,v\in\Q^\times$, write $u\preccurlyeq v$ if, for all but finitely many rational primes
$p$ and every positive integer $n$,
\[
       u^n\equiv1\pmod p\quad\Longrightarrow\quad v^n\equiv1\pmod p.
\]
Equivalently,
\[
       \ord_p(v)\mid \ord_p(u)
\]
for almost all $p$ at which both elements are defined.
\end{definition}
The equivalence follows by taking $n=\ord_p(u)$ in one direction and using divisibility of
orders in the other.

\begin{theorem}[Corrales--Rodrig\'a\~nez--Schoof]\label{thm:CRS}
Let $F$ be a number field and $u,v\in F^\times$.  If, for every integer $n$ and almost all
prime ideals $\mathfrak p$,
\[
   u^n\equiv1\pmod{\mathfrak p}
       \quad\Longrightarrow\quad
   v^n\equiv1\pmod{\mathfrak p},
\]
then $v$ is an integral power of $u$.
\end{theorem}

In the rational positive setting needed here, the conclusion has no root-of-unity ambiguity.

\begin{corollary}[Exact support reformulation]\label{cor:supportAE}
A candidate exponent is integral if either
\[
       2\preccurlyeq M
       \qquad\text{or}\qquad
       3\preccurlyeq A.
\]
In fact, $2\preccurlyeq M$ implies $M=2^k$ for some $k\in\Z$, and then $x=k$.
\end{corollary}

\subsection{Cyclotomic divisibility target}

For a prime $p\nmid2$ let $d=\ord_p(2)$.  Then $p\mid2^d-1$.  Thus the first support
condition would follow from the divisibility sequence implication
\[
       2^n-1\mid M^n-1\qquad(n\ge1).
       \tag{S2}\label{eq:S2}
\]
It is enough that \eqref{eq:S2} hold for every sufficiently large $n$: for fixed $p$, choose two
consecutive integers $k,k+1$ with both $kd,(k+1)d$ in the valid range.  Then
$M^{kd}\equiv M^{(k+1)d}\equiv1\pmod p$, hence $M^d\equiv1\pmod p$.

Condition \eqref{eq:S2} is far stronger than the original real equation, but it identifies a
concrete algebraic bridge.  It can be phrased on the $3$-smooth monoid.

\begin{definition}[Candidate monoid map]\label{def:phi}
Let
\[
       \mathcal U=\{2^a3^b:a,b\in\Z_{\ge0}\}.
\]
Define the monoid homomorphism
\[
       \Phi:\mathcal U\longrightarrow\Z_{>0},
       \qquad
       \Phi(2^a3^b)=M^aA^b.
\]
\end{definition}

\begin{definition}[Congruence preservation]
A map $f:E\to\Z$ on a subset $E\subset\Z$ is \emph{congruence-preserving} if
\[
       u-v\mid f(u)-f(v)
       \qquad(u,v\in E).
\]
\end{definition}

\begin{proposition}[Congruence-preservation criterion]\label{prop:CP}
If the candidate map $\Phi$ is congruence-preserving on $\mathcal U$, then $x\in\Z$.
The same conclusion holds if congruence preservation is known only for the pairs
$(2^n,1)$ for all sufficiently large $n$.
\end{proposition}
\begin{proof}
Since $1,2^n\in\mathcal U$ and $\Phi(1)=1$, congruence preservation gives
\[
       2^n-1\mid\Phi(2^n)-\Phi(1)=M^n-1.
\]
Now apply \cref{cor:supportAE}, with the preceding two-consecutive-multiples argument in the
eventual case.
\end{proof}

\boxedparagraph{deepgreen}{\textbf{Exact new target.}  It is unnecessary to attack the
four-exponentials statement in full generality.  It would suffice to prove the single divisibility
law $u-v\mid\Phi(u)-\Phi(v)$ on the sparse monoid of $3$-smooth integers, or even only the
one-dimensional family $2^n-1\mid M^n-1$.  Sections~\ref{sec:additive}--\ref{sec:bhargava}
show both why ordinary additive identities do not yield this law and how $p$-orderings could
encode the missing divisibility.}

\subsection{Positive-density support failures from Kummer theory}

The support theorem gives a qualitative dichotomy.  For rational integers one can also see the
failure set directly and compute its density.

\begin{theorem}[Pairwise Kummer discordance]\label{thm:pairdiscord}
Let $u,v\in\Q_{>0}^{\times}$ be multiplicatively independent.  For all but finitely many odd
primes $q$, among rational primes $p\equiv1\pmod q$ the four events
\[
\begin{array}{c|c}
\text{event}&\text{conditional density among }p\equiv1\pmod q\\ \hline
u\text{ and }v\text{ are }q\text{th powers modulo }p&1/q^2\\
u\text{ is a }q\text{th power, }v\text{ is not}&(q-1)/q^2\\
v\text{ is a }q\text{th power, }u\text{ is not}&(q-1)/q^2\\
\text{neither is a }q\text{th power}&(q-1)^2/q^2
\end{array}
\]
have the indicated Dirichlet densities.  Their unconditional densities are respectively
\[
 \frac1{(q-1)q^2},\qquad \frac1{q^2},\qquad\frac1{q^2},\qquad\frac{q-1}{q^2}.
\]
\end{theorem}
\begin{proof}
For an admissible $q$, the classes of $u$ and $v$ are independent in
$\Q(\zeta_q)^\times/(\Q(\zeta_q)^\times)^q$, so the Kummer extension
\[
   \Q(\zeta_q,u^{1/q},v^{1/q})/\Q(\zeta_q)
\]
has Galois group $\F_q^2$.  For $p\equiv1\pmod q$, the two $q$th-power residue characters are
the two coordinates of Frobenius.  Chebotarev makes this vector uniform in $\F_q^2$.  A
coordinate is zero exactly when the corresponding rational number is a $q$th power modulo
$p$.  Counting zero/nonzero coordinates gives the conditional densities.  Multiplication by
the density $1/(q-1)$ of $p\equiv1\pmod q$ gives the unconditional values.
\end{proof}

Take $u=2$ and $v=M$.  If $M$ is not a power of $2$, the pair is multiplicatively independent.
For every admissible $q$, \cref{thm:pairdiscord} supplies a positive-density set of primes such
that $2$ is a $q$th power but $M$ is not.  At each such prime, with
$n=(p-1)/q$,
\[
        2^n\equiv1\pmod p,
        \qquad
        M^n\not\equiv1\pmod p.
\]
Thus every admissible $q$ produces explicit failures of $2\preccurlyeq M$.  The reverse event
also has positive density, so the local orders of $2$ and $M$ are generically incomparable.

\begin{corollary}[Rational support theorem, density form]\label{cor:densitysupport}
For $M\in\Z_{>0}$, the following are equivalent:
\begin{enumerate}[label=(\roman*),leftmargin=2em]
\item $M$ is a nonnegative integral power of $2$;
\item $2\preccurlyeq M$;
\item for one, equivalently every, sufficiently large admissible prime $q$, the event
      ``$2$ is a $q$th power modulo $p$ but $M$ is not'' has density zero among
      $p\equiv1\pmod q$.
\end{enumerate}
\end{corollary}

This is a special rational case of the support theorem with additional quantitative information.
Its limitation is equally clear: the real identity $M=2^x$ gives no known reason for the
$q$th-power status of $M$ to follow that of $2$ when $x$ is not integral.  The next section
measures the stronger requirement that the two output residues arise from one common exponent.

\section{Prime-power residue logarithms and common-exponent coherence}
\label{sec:coherence}

\subsection{Residue-log vectors}

Fix an admissible odd prime $q$, an integer $e\ge1$, and put
\[
      Q=q^e,\qquad R_Q=\Z/Q\Z.
\]
Let $p\equiv1\pmod Q$ be a prime not dividing $6MA$.  Choose an element
$\omega_p\in\F_p^\times$ of exact order $Q$.  For $g\in\{2,3,M,A\}$ define
$\lambda_{p,Q}(g)\in R_Q$ by
\[
        g^{(p-1)/Q}=\omega_p^{\lambda_{p,Q}(g)}.
        \tag{4.1}\label{eq:reslog}
\]
Replacing $\omega_p$ by $\omega_p^u$, with $u\in R_Q^\times$, multiplies all four
coordinates by $u^{-1}$.  Consequently any condition invariant under simultaneous multiplication
by a unit is intrinsic.

Write
\[
 \bm\lambda_{p,Q}
 =\bigl(\lambda_{p,Q}(2),\lambda_{p,Q}(3),
        \lambda_{p,Q}(M),\lambda_{p,Q}(A)\bigr).
\]
Kummer reciprocity identifies this vector, up to the harmless common unit, with the Frobenius
of a prime of $K_Q$ in $L_Q/K_Q$.

\begin{lemma}[Uniform residue-log distribution]\label{lem:uniform}
Let $E\subseteq\Gal(L_Q/K_Q)$ be invariant under multiplication of all coordinates by units in
$R_Q^\times$.  Then the set of rational primes $p\equiv1\pmod Q$ for which
$\bm\lambda_{p,Q}\in E$ has conditional Dirichlet density
\[
       \frac{|E|}{|\Gal(L_Q/K_Q)|}.
\]
\end{lemma}
\begin{proof}
Chebotarev makes Frobenius elements of prime ideals of $K_Q$ uniform in the abelian group
$\Gal(L_Q/K_Q)$.  A rational prime $p\equiv1\pmod Q$ splits completely in $K_Q$.  Its
$\varphi(Q)$ primes correspond to the choices of the image of $\zeta_Q$, hence to replacing
$\omega_p$ by $\omega_p^u$ for $u\in R_Q^\times$.  Since $E$ is invariant under this common
rescaling, either all primes of $K_Q$ above $p$ satisfy the condition or none do.  The prime-ideal
Chebotarev density therefore descends unchanged to the conditional density of rational primes.
\end{proof}

\subsection{The coherence condition}

\begin{definition}[$Q$-coherence]\label{def:coherent}
A prime $p\equiv1\pmod Q$ is \emph{$Q$-coherent} for the quadruple $(2,3,M,A)$ if there
exists $k\in R_Q$ such that
\[
\begin{split}
  \lambda_{p,Q}(M)&=k\lambda_{p,Q}(2),\\
  \lambda_{p,Q}(A)&=k\lambda_{p,Q}(3).
\end{split}
\tag{4.2}\label{eq:coherent}
\]
\end{definition}
This condition is invariant under the choice of $\omega_p$.

If there is an actual integer $k$ satisfying
$M\equiv2^k\pmod p$ and $A\equiv3^k\pmod p$, then \eqref{eq:coherent} follows after raising
to $(p-1)/Q$.  The converse need not hold: $Q$-coherence sees only a quotient of the full
discrete logarithm group.  It is therefore a deliberately coarse but tractable necessary
condition for a common finite-field exponent.

If $x=k\in\Z$, the equalities $M=2^k$, $A=3^k$ hold over $\Z$, so every prime
$p\nmid6MA$ is $Q$-coherent for every $Q\mid p-1$.  The nonintegral branches behave in the
opposite way.

\subsection{Exact count in multiplicative rank four}

\begin{theorem}[Rank-four coherence density]\label{thm:rank4density}
Assume $\rank\Gamma=4$.  Let $q$ be admissible and $Q=q^e$.  Define
\[
 N_{4}(q,e)
 =1+\frac{q(q^2-1)(q^{3e}-1)}{q^3-1}.
 \tag{4.3}\label{eq:N4}
\]
Then the conditional Dirichlet density of $Q$-coherent primes among
$p\equiv1\pmod Q$ is
\[
       P_4(q,e)=\frac{N_4(q,e)}{q^{4e}}.
       \tag{4.4}\label{eq:P4}
\]
The unconditional density is
\[
       \frac{N_4(q,e)}{\varphi(Q)q^{4e}}.
\]
Moreover,
\[
       P_4(q,e)
       \sim \frac{q(q+1)}{q^2+q+1}\,q^{-e}
       \qquad(e\to\infty).
       \tag{4.5}\label{eq:P4asym}
\]
\end{theorem}
\begin{proof}
By \cref{lem:kummermodule,lem:uniform}, the residue-log vector is uniform in
$R_Q^4$.  Write its base and output halves as
\[
       \bm b=(u,v),\qquad \bm y=(m,a).
\]
For fixed $\bm b$, the coherent outputs are exactly the cyclic submodule
$R_Q\bm b$.  If $\bm b=0$, this submodule has one element.  Otherwise define its content
$s$ by
\[
        s=\min\{v_q(u),v_q(v),e\},\qquad 0\le s<e.
\]
Put $t=e-s$.  There are
\[
        q^{2t}-q^{2t-2}
\]
base vectors of content $s$, and the cyclic submodule they generate has $q^t$ elements.  Hence
the number of coherent quadruples is
\[
\begin{split}
  N_4(q,e)
   &=1+\sum_{t=1}^{e}(q^{2t}-q^{2t-2})q^t\\
   &=1+\frac{q(q^2-1)(q^{3e}-1)}{q^3-1}.
\end{split}
\]
Division by $|R_Q^4|=q^{4e}$ and application of \cref{lem:uniform} gives the conditional
density.  The primes $p\equiv1\pmod Q$ have density $1/\varphi(Q)$, giving the
unconditional value.  The leading term in \eqref{eq:N4} yields \eqref{eq:P4asym}.
\end{proof}

At depth one,
\[
        P_4(q,1)=\frac{q^3-q+1}{q^4}
        =\frac1q-\frac1{q^3}+\frac1{q^4}.
        \tag{4.6}
\]

\subsection{Exact count in multiplicative rank three}

The rank-three relation cuts the ambient Kummer group by one linear equation, but the coherence
variety also intersects that hyperplane nontransversely at imprimitive base vectors.  The exact
prime-power count therefore differs from simply multiplying the rank-four probability by $q$.

\begin{theorem}[Rank-three coherence density]\label{thm:rank3density}
Assume $\rank\Gamma=3$, let $\bm c$ be the primitive relation, and let $q$ be admissible in
the sense of \cref{def:admissible}.  Put $Q=q^e$ and
\[
 N_3(q,e)
 =1+\frac{q(q^{2e}-1)}{q+1}.
 \tag{4.7}\label{eq:N3}
\]
Then the conditional Dirichlet density of $Q$-coherent primes among
$p\equiv1\pmod Q$ is
\[
       P_3(q,e)=\frac{N_3(q,e)}{q^{3e}}.
       \tag{4.8}\label{eq:P3}
\]
The unconditional density is
\[
       \frac{N_3(q,e)}{\varphi(Q)q^{3e}},
\]
and
\[
       P_3(q,e)\sim\frac{q}{q+1}\,q^{-e}
       \qquad(e\to\infty).
       \tag{4.9}\label{eq:P3asym}
\]
\end{theorem}
\begin{proof}
The Kummer group is the hyperplane
\[
 H_e=\{(u,v,m,a)\in R_Q^4:
        c_1u+c_2v+c_3m+c_4a=0\},
\]
of cardinality $q^{3e}$.  A coherent vector has the form
$(u,v,ku,kv)$.  Thus it lies in $H_e$ precisely when
\[
       (c_1+c_3k)u+(c_2+c_4k)v=0\pmod Q.
       \tag{4.10}\label{eq:rank3eq}
\]
The zero base vector contributes one coherent vector.

Now let $(u,v)$ have content $s<e$ and put $t=e-s$.  Divide by $q^s$ and work modulo
$q^t$ with a primitive vector $(u',v')$.  Set
\[
   \alpha=c_1u'+c_2v',\qquad
   \delta=c_3u'+c_4v'.
\]
Because $q\nmid\Delta(\bm c)$, the matrix
\[
      \begin{pmatrix}c_1&c_2\\c_3&c_4\end{pmatrix}
\]
is invertible modulo $q^t$.  Therefore $(u',v')\mapsto(\alpha,\delta)$ permutes the
primitive vectors of $(\Z/q^t\Z)^2$.  Equation \eqref{eq:rank3eq} becomes
\[
        \alpha+k\delta=0\pmod{q^t}.
\]
For a primitive pair $(\alpha,\delta)$ this has a solution exactly when $\delta$ is a unit,
and then the solution $k$ modulo $q^t$ is unique.  The number of such pairs is
\[
       q^t(q^t-q^{t-1})=q^{2t}-q^{2t-1}.
\]
Each gives one distinct coherent output.  Summing over contents,
\[
\begin{split}
 N_3(q,e)
   &=1+\sum_{t=1}^{e}(q^{2t}-q^{2t-1})\\
   &=1+\frac{q(q^{2e}-1)}{q+1}.
\end{split}
\]
Divide by $|H_e|=q^{3e}$ and invoke \cref{lem:uniform}.  The remaining assertions follow
as in the rank-four case.
\end{proof}

At depth one,
\[
        P_3(q,1)=\frac{q^2-q+1}{q^3}
        =\frac1q-\frac1{q^2}+\frac1{q^3}.
        \tag{4.11}
\]
For every $q>1$,
$P_4(q,1)-P_3(q,1)=(q-1)^2/q^4>0$; the rank-four branch is slightly more coherent because
there is no relation hyperplane to constrain the output.

\subsection{The density gap}

\begin{theorem}[Finite-place density gap]\label{thm:densitygap}
Let $(M,A)$ arise from a candidate exponent.
\begin{enumerate}[label=(\alph*),leftmargin=2em]
\item If $x\in\Z$, then for every $Q$ and every prime $p\equiv1\pmod Q$ with
      $p\nmid6MA$, the prime $p$ is $Q$-coherent.
\item If $x\notin\Z$, then $\rank\Gamma\in\{3,4\}$, and for every admissible $q$ and
      every $e\ge1$ the conditional density of $q^e$-coherent primes is exactly
      $P_3(q,e)$ or $P_4(q,e)$ according to the rank branch.  In particular it is
      $O_q(q^{-e})$, while the incoherent density tends to $1$ as $e\to\infty$.
\end{enumerate}
\end{theorem}
\begin{proof}
Part (a) follows by taking the common scalar $k=x\bmod Q$ in
\eqref{eq:coherent}.  Part (b) combines \cref{prop:rankdichotomy,thm:rank4density,thm:rank3density}.
\end{proof}

\begin{corollary}[Quantitative sufficient criterion]\label{cor:quantcriterion}
Fix an admissible $q$ and $e\ge1$.  If one can prove, from the candidate equations, that the
conditional density of $q^e$-coherent primes is strictly larger than both
$P_3(q,e)$ and $P_4(q,e)$, then $x\in\Z$.
\end{corollary}

The criterion is deliberately quantitative.  A statement such as ``there are infinitely many
coherent primes'' is useless: both nonintegral rank branches have positive density at every fixed
depth.  One needs a density excess, coherence at almost every admissible prime, or a mechanism
that remains coherent through unbounded depth.

\boxedparagraph{deepblue}{\textbf{Interpretation.}  A hypothetical counterexample would satisfy
the exact real determinant
$\log2\,\log A-\log3\,\log M=0$, yet at depth $q^e$ its corresponding finite residue vectors
would fail to share a scalar at a proportion tending to $1$.  The conjecture asks for a bridge
showing that this extreme archimedean/finite-place mismatch cannot occur for the special integer
quadruple.  Chebotarev quantifies the mismatch perfectly; it does not remove it.}

\section{Recovering the rank branch from finite residue patterns}
\label{sec:detector}

The depth-one Kummer data contains more than the aggregate coherence probability.  In rank four,
the four $q$th-power statuses are independent.  In rank three, the unique multiplicative relation
forbids entire status patterns.  This gives a finite-place detector for the relation support.

Fix an admissible prime $q$ and a prime $p\equiv1\pmod q$.  A coordinate
$\lambda_{p,q}(g)$ is zero exactly when $g$ is a $q$th power modulo $p$.  Let
\[
       J_q=\{i\in\{1,2,3,4\}:c_i\not\equiv0\pmod q\}
\]
be the active support of the rank-three relation modulo $q$.

\begin{lemma}[Nonzero solutions of one linear equation]\label{lem:Ns}
For $s\ge0$, let $N_s(q)$ be the number of $(z_1,\dots,z_s)\in(\F_q^\times)^s$ satisfying
\[
       z_1+\cdots+z_s=0.
\]
Then
\[
       N_s(q)=\frac{(q-1)^s+(q-1)(-1)^s}{q}.
       \tag{5.1}\label{eq:Ns}
\]
In particular $N_0(q)=1$ and $N_1(q)=0$.
\end{lemma}
\begin{proof}
Using additive-character orthogonality,
\[
 N_s(q)=\frac1q\sum_{t\in\F_q}
       \left(\sum_{z\in\F_q^\times}\exp\frac{2\pi i\,\mathrm{Tr}(tz)}q\right)^s.
\]
The inner sum is $q-1$ for $t=0$ and $-1$ for each $t\ne0$, yielding
\eqref{eq:Ns}.  Equivalently, one may prove the same formula by the recurrence
$N_s=(q-1)^{s-1}-N_{s-1}$.
\end{proof}

\begin{theorem}[Rank-three status-pattern law]\label{thm:patternlaw}
Assume $\rank\Gamma=3$ and let $q$ be admissible.  Prescribe, for each of the four generators,
whether its residue-log coordinate is zero or nonzero.  Let
\begin{itemize}[leftmargin=2em,itemsep=1pt]
\item $s$ be the number of active coordinates in $J_q$ prescribed to be nonzero;
\item $h$ be the number of inactive coordinates outside $J_q$ prescribed to be nonzero.
\end{itemize}
Then the number of points in the Kummer hyperplane realizing the pattern is
\[
       N_s(q)(q-1)^h.
       \tag{5.2}\label{eq:patterncount}
\]
Its conditional density among primes $p\equiv1\pmod q$ is
\[
       \frac{N_s(q)(q-1)^h}{q^3}.
       \tag{5.3}\label{eq:patterndensity}
\]
In particular, every pattern with exactly one active nonzero coordinate is forbidden.
\end{theorem}
\begin{proof}
The residue vector is uniform on the hyperplane $\bm c\cdot\bm t=0$ by
\cref{lem:uniform}.  After independently scaling each active coordinate by the nonzero
coefficient $c_i$, the relation is a sum of the $s$ prescribed nonzero active variables.  There
are $N_s(q)$ choices by \cref{lem:Ns}.  Inactive nonzero coordinates are unrestricted and
contribute $(q-1)^h$.  The hyperplane has $q^3$ points.
\end{proof}

\begin{proposition}[Rank-four status law]\label{prop:rank4patterns}
If $\rank\Gamma=4$, every zero/nonzero pattern occurs.  A pattern with $z$ zero coordinates
has conditional density
\[
       \frac{(q-1)^{4-z}}{q^4}.
\]
Thus the number of zero coordinates is binomial with parameters $4$ and $1/q$.
\end{proposition}

The forbidden-pattern law is stronger than a generic rank computation: by varying $q$, observed
forbidden patterns reveal which coefficients $c_i$ vanish modulo $q$, and the frequencies of the
remaining patterns constrain the relation support.  With full residue-log values rather than
only zero/nonzero statuses, one can recover the projective relation $[\bm c]\in\mathbb P^3(\F_q)$
from the orthogonal hyperplane.  Chinese remaindering over several admissible $q$ then reconstructs
$\bm c$ once a height bound is available.

\begin{example}[A clean rank-three test tuple]
The quadruple $(2,3,7,42)$ has relation
\[
       2\cdot3\cdot7=42,
\]
so $\bm c=(1,1,1,-1)$ and every coordinate is active for $q=5$.  The four patterns with exactly
one non-$5$th-power coordinate are therefore forbidden.  In the computation of
\cref{sec:experiments}, each occurs zero times among $87{,}062$ primes $p\le5{,}000{,}000$ with
$p\equiv1\pmod5$.
\end{example}

\section{Exact computational checks}
\label{sec:experiments}

All experiments in this section use exact integer modular arithmetic.  They are included for
three reasons: to catch count-normalization errors, to test the descent from prime ideals to
rational primes, and to provide small local witnesses suitable for later formalization.  They do
not test whether a candidate counterexample exists.

\subsection{Direct finite-ring enumeration}

Before sampling primes, the two numerator formulas were checked by exhaustive enumeration of
all quadruples over $\Z/q^e\Z$.  In the rank-three column the hyperplane was
$u+v+m-a=0$, corresponding to the test relation $2\cdot3\cdot7=42$.
\begin{center}
\begin{tabular}{rrrrr}
\toprule
$q$ & $e$ & exhaustive $N_4$ & formula $N_4(q,e)$ & exhaustive/formula $N_3$\\
\midrule
$3$ & $1$ & $25$    & $25$    & $7/7$\\
$3$ & $2$ & $673$   & $673$   & $61/61$\\
$5$ & $1$ & $121$   & $121$   & $21/21$\\
$5$ & $2$ & $15{,}121$ & $15{,}121$ & $521/521$\\
\bottomrule
\end{tabular}
\end{center}
This check is independent of prime generation and directly tests the zero-divisor cases at
$e=2$.

\subsection{Pairwise discordance}

Take $(u,v)=(2,14)$ and $q=5$.  Neither generator is divisible by $5$, and their classes are
independent.  Among primes $p\le5{,}000{,}000$ with $p\equiv1\pmod5$, the observed counts are:
\begin{center}
\begin{tabular}{lrrr}
\toprule
Status modulo $p$ & Count & Frequency & Predicted frequency\\
\midrule
both $5$th powers & $3{,}339$ & $0.03835198$ & $0.04000000$\\
$2$ a $5$th power, $14$ not & $14{,}006$ & $0.16087386$ & $0.16000000$\\
$14$ a $5$th power, $2$ not & $13{,}980$ & $0.16057522$ & $0.16000000$\\
neither a $5$th power & $55{,}737$ & $0.64019894$ & $0.64000000$\\
\midrule
Total & $87{,}062$ & $1$ & $1$\\
\bottomrule
\end{tabular}
\end{center}

For several values of $q$, the least prime witnesses found below the same limit are shown next.
Here $n=(p-1)/q$; the first witness has $2^n\equiv1$ and $14^n\not\equiv1$, while the second
has the reverse status.
\begin{center}
\begin{tabular}{rrrrr}
\toprule
$q$ & first $p$ & first $n$ & reverse $p$ & reverse $n$\\
\midrule
3  & 31   & 10  & 13   & 4\\
5  & 151  & 30  & 41   & 8\\
7  & 631  & 90  & 71   & 10\\
11 & 331  & 30  & 397  & 36\\
13 & 4421 & 340 & 547  & 42\\
17 & 1429 & 84  & 2347 & 138\\
19 & 1103 & 58  & 1597 & 84\\
\bottomrule
\end{tabular}
\end{center}
These are concrete finite certificates of the two support failures.

\subsection{Common-exponent coherence at depths \texorpdfstring{$5$ and $25$}{5 and 25}}

Two toy quadruples isolate the two rank branches:
\[
       (2,3,7,11)\quad\text{has rank }4,
\]
while
\[
       (2,3,7,42)\quad\text{has rank }3,
       \qquad 2\cdot3\cdot7=42.
\]
The next table compares the observed coherence frequencies with
\cref{thm:rank4density,thm:rank3density}.  The $Q=5$ rows use primes at most $5{,}000{,}000$;
the $Q=25$ rows use primes at most $10{,}000{,}000$.  The $z$ column is the deviation divided
by the binomial standard deviation $\sqrt{NP(1-P)}$.

\begin{center}
\begin{tabular}{llrrrrr}
\toprule
Rank & $Q$ & Eligible primes & Coherent & Observed & Predicted & $z$\\
\midrule
4 & $5$  & $87{,}061$ & $16{,}760$ & $0.19250870$ & $0.19360000$ & $-0.81$\\
3 & $5$  & $87{,}062$ & $14{,}476$ & $0.16627231$ & $0.16800000$ & $-1.36$\\
4 & $25$ & $33{,}155$ & $1{,}281$  & $0.03863671$ & $0.03870976$ & $-0.07$\\
3 & $25$ & $33{,}155$ & $1{,}098$  & $0.03311718$ & $0.03334400$ & $-0.23$\\
\bottomrule
\end{tabular}
\end{center}
The depth-$25$ predicted numerators are
\[
       N_4(5,2)=15{,}121,
       \qquad
       N_3(5,2)=521.
\]
The agreement is especially useful because prime-power rings contain zero divisors; a mistaken
field-only count would generally first appear at $e=2$.

\subsection{Status-pattern checks}

For the rank-four tuple at $q=5$, grouping by the number of coordinates that are $5$th powers
gives:
\begin{center}
\begin{tabular}{rrrr}
\toprule
Number of $5$th-power coordinates & Count & Observed & Predicted\\
\midrule
0 & $35{,}627$ & $0.40921882$ & $0.40960000$\\
1 & $35{,}800$ & $0.41120594$ & $0.40960000$\\
2 & $13{,}354$ & $0.15338671$ & $0.15360000$\\
3 & $2{,}143$  & $0.02461493$ & $0.02560000$\\
4 & $137$      & $0.00157361$ & $0.00160000$\\
\bottomrule
\end{tabular}
\end{center}
For the rank-three tuple, all four one-nonresidue patterns
\[
1110,\quad 1101,\quad 1011,\quad 0111
\]
(where $1$ means ``$5$th power'') had count zero, exactly as
\cref{thm:patternlaw} predicts.  Selected nonzero patterns were:
\begin{center}
\begin{tabular}{rrrr}
\toprule
Pattern & Count & Observed & Predicted\\
\midrule
0000 & $36{,}185$ & $0.41562335$ & $0.41600000$\\
0001 & $8{,}387$  & $0.09633365$ & $0.09600000$\\
0011 & $2{,}781$  & $0.03194275$ & $0.03200000$\\
1111 & $649$      & $0.00745446$ & $0.00800000$\\
\bottomrule
\end{tabular}
\end{center}
The maximum absolute frequency error over all sixteen patterns was $0.001103$.

\subsection{What the computations do and do not establish}

The tests strongly validate the finite counts and implementation.  They do not add an assumption
such as GRH, because the theoretical densities are unconditional consequences of Chebotarev.
They also do not provide statistical evidence for the Alaoglu--Erd\H{o}s conjecture: the toy
tuples were chosen solely to have known multiplicative rank.  Their role is proof engineering.

\section{The additive skeleton of the \texorpdfstring{$3$-smooth}{3-smooth} monoid}
\label{sec:additive}

The support criterion in \cref{prop:CP} asks for congruence preservation on
$\mathcal U=\{2^a3^b\}$.  A natural first attempt is to exploit additive identities in
$\mathcal U$ and transport them through $\Phi$.  This section classifies all such identities
and proves that the attempt collapses to one fixed modulus.

\subsection{Classification of primitive additive triples}

Call a solution $u+v=w$ in $\mathcal U$ \emph{primitive} if
$\gcd(u,v,w)=1$.

\begin{theorem}[Complete primitive additive skeleton]\label{thm:additiveclassification}
Up to interchanging $u$ and $v$, the primitive solutions of
\[
       u+v=w,\qquad u,v,w\in\mathcal U,
\]
are exactly
\[
       1+1=2,\qquad
       1+2=3,\qquad
       1+3=4,\qquad
       1+8=9.
       \tag{7.1}\label{eq:fourtriples}
\]
Every nonprimitive solution is obtained by multiplying one of these four by an element of
$\mathcal U$.
\end{theorem}
\begin{proof}
In a primitive solution the three numbers are pairwise coprime, since
$\gcd(u,w)=\gcd(u,v)$ and similarly for $v$.  If both $u$ and $v$ are greater than one, their
prime supports must be disjoint, so up to order $u=2^a$, $v=3^b$ with $a,b>0$.  But then $w$
is coprime to both $2$ and $3$, forcing $w=1$, impossible.  Thus one summand is $1$.

The case $1+1=2$ is immediate.  The remaining equations are
\[
       3^b-2^a=1
       \qquad\text{or}\qquad
       2^a-3^b=1.
\]
For the first, $a=1$ gives $(a,b)=(1,1)$.  If $a\ge2$, reduction modulo $4$ makes $b$ even,
say $b=2m$.  Then
\[
       (3^m-1)(3^m+1)=2^a.
\]
Both factors are powers of $2$, have gcd $2$, and differ by $2$; hence they are $2$ and $4$.
Thus $m=1$, giving $(a,b)=(3,2)$.

For the second equation, reduction modulo $3$ makes $a$ even, say $a=2m$.  Then
\[
       (2^m-1)(2^m+1)=3^b.
\]
The coprime factors are powers of $3$ differing by $2$, so they are $1$ and $3$.  Hence
$m=1$ and $(a,b)=(2,1)$.  This proves \eqref{eq:fourtriples}.  Dividing any nonprimitive
solution by its common gcd, which is again $3$-smooth, gives the final assertion.
\end{proof}

No use of Catalan--Mih\u ailescu is needed; the two special exponential equations factor after a
single congruence observation.

\subsection{Defects of the four primitive relations}

Transport the four identities through $\Phi$.  With signs chosen for convenience, define
\[
\begin{aligned}
 d_{11}&=M-2,\\
 d_{12}&=A-M-1,\\
 d_{13}&=M^2-A-1,\\
 d_{18}&=A^2-M^3-1.
\end{aligned}
\tag{7.2}\label{eq:defects}
\]
They satisfy the exact identities
\[
       d_{13}=(M+1)d_{11}-d_{12},
       \tag{7.3}\label{eq:d13}
\]
and
\[
       d_{18}=-M(M+1)d_{11}+\bigl(2(M+1)+d_{12}\bigr)d_{12}.
       \tag{7.4}\label{eq:d18}
\]
Moreover
\[
       d_{12}=(A-3)-(M-2).
\]
Thus all four primitive defects generate the same ideal as $M-2$ and $A-3$.

\begin{theorem}[Additive defect ideal]\label{thm:defectideal}
Let $\mathcal D$ be the ideal of $\Z$ generated by all transported additive defects
\[
       \Phi(u)+\Phi(v)-\Phi(w)
\]
with $u+v=w$ in $\mathcal U$.  Then
\[
       \mathcal D=(M-2,A-3)
       =C\Z,
       \qquad C:=\gcd(M-2,A-3).
       \tag{7.5}\label{eq:C}
\]
\end{theorem}
\begin{proof}
By \cref{thm:additiveclassification}, every additive triple is a $3$-smooth multiple of one
of the four primitive triples.  Under $\Phi$, multiplication by $2^a3^b$ multiplies the defect
by $M^aA^b$.  Hence all defects lie in the ideal generated by the four numbers in
\eqref{eq:defects}.  Equations \eqref{eq:d13}--\eqref{eq:d18} reduce that ideal to
$(d_{11},d_{12})$, which equals $(M-2,A-3)$.  Conversely, $d_{11}$ and $d_{12}$ themselves
come from the first two primitive identities, so equality holds.
\end{proof}

The next theorem shows that allowing arbitrary signed polynomial combinations of additive
identities creates no new modulus.

\begin{theorem}[Universal polynomial collapse]\label{thm:polycollapse}
For every $P\in\Z[X,Y]$ with $P(2,3)=0$,
\[
       C\mid P(M,A).
       \tag{7.6}\label{eq:polydiv}
\]
Conversely, the ideal generated by all integers $P(M,A)$ with $P(2,3)=0$ is exactly $C\Z$.
Equivalently, a positive integer $D$ has the property
\[
       \Phi(u)\equiv u\pmod D
       \qquad\text{for every }u\in\mathcal U.
       \tag{7.7}\label{eq:largestmod}
\]
if and only if $D\mid C$.  Hence $C$ is the largest such modulus when $C>0$; when $C=0$
every positive modulus works.
\end{theorem}
\begin{proof}
The evaluation kernel at $(2,3)$ is
\[
       \ker\bigl(\Z[X,Y]\xrightarrow{P\mapsto P(2,3)}\Z\bigr)
       =(X-2,Y-3).
\]
Thus $P=(X-2)U+(Y-3)V$ for some $U,V\in\Z[X,Y]$, and
$P(M,A)$ is divisible by $C$.  Conversely, $P=X-2$ and $P=Y-3$ show that the generated
ideal contains $(M-2,A-3)=C\Z$.

If $D\mid C$, then $M\equiv2$ and $A\equiv3\pmod D$, so every monomial satisfies
$M^aA^b\equiv2^a3^b\pmod D$.  Conversely, any positive modulus with that property must divide
both $M-2$ and $A-3$, hence divide $C$.
\end{proof}

\begin{warning}
The conclusion is not that additive information is useless.  It is that every identity obtained
solely by evaluating an integral polynomial relation true at $(2,3)$ can see only the fixed
anchor $C$.  A proof of congruence preservation must use varying differences $u-v$, multiplicative
orders, or prime-dependent geometry; it cannot arise from accumulating more polynomial identities
at the single point $(2,3)$.
\end{warning}

\section{Dynamic anchors and a subexponential no-go theorem}
\label{sec:dynamic}

The fixed modulus $C$ of \cref{thm:polycollapse} compares $(M,A)$ with $(2,3)$ at one time.
A stronger attempt compares the $n$th powers and allows the modulus to vary:
\[
       C_n:=\gcd(M^n-2^n,A^n-3^n).
       \tag{8.1}\label{eq:Cn}
\]
For a nonintegral candidate neither difference is identically zero.

\begin{proposition}[Dynamic polynomial collapse]\label{prop:dynamiccollapse}
For each $n\ge1$, the ideal generated by
\[
       \{P(M^n,A^n):P\in\Z[X,Y],\ P(2^n,3^n)=0\}
\]
is exactly $C_n\Z$.  Equivalently, a positive integer $D$ has
\[
       M^{an}A^{bn}\equiv2^{an}3^{bn}\pmod D
       \qquad(a,b\ge0)
\]
if and only if $D\mid C_n$.  Thus $C_n$ is the largest such modulus when $C_n>0$; if
$C_n=0$, every positive modulus works.
Moreover, if $m\mid n$, then $C_m\mid C_n$.
\end{proposition}
\begin{proof}
Apply the proof of \cref{thm:polycollapse} with the evaluation point
$(2^n,3^n)$.  Its kernel is $(X-2^n,Y-3^n)$, and evaluation at $(M^n,A^n)$ produces the
ideal $(M^n-2^n,A^n-3^n)=C_n\Z$.  If $m\mid n$, then
$X^m-Y^m$ divides $X^n-Y^n$, so both generators of the $m$th ideal divide their
$n$th counterparts.
\end{proof}

For a prime $p\nmid6MA$,
\[
 p\mid C_n
 \quad\Longleftrightarrow\quad
 \left(\frac M2\right)^n\equiv1\pmod p
 \ \text{and}\ 
 \left(\frac A3\right)^n\equiv1\pmod p.
 \tag{8.2}\label{eq:Cnorders}
\]
Thus $C_n$ records simultaneous return times of the two rational ratios $M/2$ and $A/3$.

\begin{lemma}[Independence of the dynamic ratios]\label{lem:ratioind}
If $x\ne1$, the rational numbers $M/2$ and $A/3$ are multiplicatively independent.
\end{lemma}
\begin{proof}
A relation $(M/2)^r(A/3)^s=1$ gives, after substituting $M=2^x$, $A=3^x$ and taking
real logarithms,
\[
       (x-1)(r\log2+s\log3)=0.
\]
Since $x\ne1$ and $2,3$ are multiplicatively independent, $r=s=0$.
\end{proof}

\begin{theorem}[Subexponential dynamic anchors]\label{thm:Cnsubexp}
For a nonintegral candidate,
\[
       \log C_n=o(n)
       \qquad(n\to\infty).
       \tag{8.3}\label{eq:Cnsubexp}
\]
Equivalently, for every $\varepsilon>0$, one has $C_n<e^{\varepsilon n}$ for all sufficiently
large $n$.
\end{theorem}
\begin{proof}
Put $U=M/2$ and $V=A/3$.  By \cref{lem:ratioind}, $U$ and $V$ are multiplicatively
independent rational $S$-units for the fixed finite set of primes dividing $6MA$.  The
Corvaja--Zannier $S$-unit gcd estimate, extending the Bugeaud--Corvaja--Zannier theorem, gives
an $o(n)$ upper bound for the contribution outside $S$ to the generalized gcd of
$U^n-1$ and $V^n-1$.

It remains to control primes in $S$.  If a prime $p\notin\{2,3\}$ divides $M$, then
$v_p(M^n-2^n)=0$; if it divides $A$, then $v_p(A^n-3^n)=0$.  Thus such primes contribute
nothing to the common gcd unless both relevant terms are $p$-adic units, in which case the
standard lifting-the-exponent argument gives $O(v_p(n)+1)=O(\log n)$.  At $p=2$, either
$A$ is even and $A^n-3^n$ is odd, or $A$ is odd and the same lifting argument gives
$v_2(A^n-3^n)=O(\log n)$.  The prime $3$ is symmetric, using $M^n-2^n$.  Hence the total
contribution from the fixed set $S$ is $O(\log n)$.  Combining it with the outside-$S$
estimate proves \eqref{eq:Cnsubexp}.
\end{proof}

\begin{remark}
The theorem also holds for integral exponents $x=k>1$; it is a limitation theorem for the
anchor method, not an integer/noninteger discriminator.  The exceptional exponent $x=1$ makes
both differences zero.
\end{remark}

\subsection{Why the additive route stops}

A support proof along \eqref{eq:S2} would have to transfer all prime divisors of $2^n-1$ to
$M^n-1$.  The modulus $2^n-1$ has exponential logarithmic size.  By contrast, every modulus
produced by simultaneous polynomial return of $(M^n,A^n)$ to $(2^n,3^n)$ divides $C_n$, whose
logarithm is $o(n)$.  Thus there is an exponential gap:
\[
      \log(2^n-1)\sim n\log2,
      \qquad
      \log C_n=o(n).
\]
This does not refute all uses of addition.  It refutes the specific strategy of accumulating
ordinary polynomial identities at the two evaluation points and hoping their common defect
becomes large enough to enforce support.  Any successful additive argument must retain
prime-by-prime order information that the single gcd $C_n$ discards.

\section{Local congruence kernels and the Bhargava-ordering route}
\label{sec:bhargava}

This section converts global congruence preservation into exact finite cyclic-group statements.
It then explains how $p$-orderings offer a disciplined way to attack those statements without
claiming that the existing theory solves them automatically.

\subsection{Kernel inclusion on exponent lattices}

For an odd prime $p\nmid6MA$ and $m\ge1$, define homomorphisms
\[
\begin{aligned}
 \rho_{p^m}:\Z^2&\longrightarrow(\Z/p^m\Z)^\times,
 & (a,b)&\longmapsto2^a3^b,\\
 \rho'_{p^m}:\Z^2&\longrightarrow(\Z/p^m\Z)^\times,
 & (a,b)&\longmapsto M^aA^b.
\end{aligned}
\tag{9.1}\label{eq:rhomaps}
\]
Although $\Phi$ was defined only for nonnegative exponents, any difference in $\Z^2$ can be
realized by adding a sufficiently large common vector to two nonnegative exponent pairs.
Consequently the local congruence condition on $\mathcal U$ is exactly a kernel inclusion.

\begin{proposition}[Local kernel criterion]\label{prop:kernelcriterion}
For fixed $p\nmid6MA$ and $m\ge1$, the following are equivalent.
\begin{enumerate}[label=(\roman*),leftmargin=2em]
\item Whenever $u,v\in\mathcal U$ satisfy $u\equiv v\pmod{p^m}$, one has
      $\Phi(u)\equiv\Phi(v)\pmod{p^m}$.
\item $\ker\rho_{p^m}\subseteq\ker\rho'_{p^m}$.
\item The rule
\[
      2^a3^b\longmapsto M^aA^b
\]
defines a well-defined group homomorphism
\[
      f_{p^m}:\langle2,3\rangle\longrightarrow(\Z/p^m\Z)^\times.
\]
\end{enumerate}
\end{proposition}
\begin{proof}
The equivalence of (ii) and (iii) is the universal property of a quotient.  To compare with (i),
write $u=2^a3^b$ and $v=2^c3^d$.  Since all four generators are units modulo $p^m$,
$u\equiv v$ is equivalent to $(a-c,b-d)\in\ker\rho_{p^m}$, and the corresponding output
congruence is membership in $\ker\rho'_{p^m}$.  Negative coordinate differences cause no issue,
as both exponent pairs may be shifted into $\Z_{\ge0}^2$.
\end{proof}

For odd $p$, the unit group $(\Z/p^m\Z)^\times$ is cyclic.  This makes every local monoid
homomorphism a common power map.

\begin{theorem}[Local common-exponent criterion]\label{thm:localcommon}
Let $p$ be an odd prime with $p\nmid6MA$, and let $m\ge1$.  The equivalent conditions of
\cref{prop:kernelcriterion} hold if and only if there exists an integer $k_{p,m}$ such that
\[
       M\equiv2^{k_{p,m}}\pmod{p^m},
       \qquad
       A\equiv3^{k_{p,m}}\pmod{p^m}.
       \tag{9.2}\label{eq:fullcommon}
\]
\end{theorem}
\begin{proof}
If \eqref{eq:fullcommon} holds, raising congruent elements of $\mathcal U$ to the common power
$k_{p,m}$ proves local congruence preservation.

Conversely, let $H=\langle2,3\rangle$ in the cyclic group
$G=(\Z/p^m\Z)^\times$, and let $f:H\to G$ be the homomorphism from
\cref{prop:kernelcriterion}.  The image $f(H)$ has order dividing $|H|$.  A cyclic group has a
unique subgroup of each order dividing $|G|$, so every element whose order divides $|H|$ lies
in $H$; hence $f(H)\subseteq H$.  Thus $f$ is an endomorphism of the cyclic group $H$, and every
such endomorphism is $h\mapsto h^k$ for some integer $k$.  Evaluating at $2$ and $3$ gives
\eqref{eq:fullcommon}.
\end{proof}

\begin{corollary}[Global congruence preservation as local exponents]\label{cor:globalCP}
Away from the finite set of primes dividing $6MA$, the map $\Phi$ is congruence-preserving if
and only if compatible common exponents $k_{p,m}$ as in \eqref{eq:fullcommon} exist for every
$p$ and $m$.  For the support conclusion alone, it is enough to have a common exponent modulo
$p$ for almost all $p$.
\end{corollary}

The depth-$Q$ coherence of \cref{def:coherent} is a quotient shadow of
\eqref{eq:fullcommon}: a full common exponent modulo $p$ implies $Q$-coherence for every
$Q\mid p-1$.  The exact densities therefore quantify how rarely the first necessary quotient
condition holds in the nonintegral rank branches.

\subsection{Bhargava \texorpdfstring{$p$-orderings}{p-orderings} as a concrete coordinate system}

Bhargava introduced $p$-orderings of an arbitrary subset of a Dedekind domain to describe
integer-valued polynomial functions and generalized factorials.  For a subset
$E\subset\Z$ and a prime $p$, a $p$-ordering is a sequence
$e_0,e_1,\ldots$ chosen recursively so that $e_n$ minimizes
\[
       v_p\!\left(\prod_{j<n}(e-e_j)\right)
\]
among $e\in E$.  The resulting valuation sequence is independent of the choices and defines
the generalized factorial of $E$.

For $E=\mathcal U$, the residues are generated multiplicatively by $2$ and $3$.  Modulo
$p^m$, the geometry of $E$ is therefore controlled by the lattice
$\ker\rho_{p^m}\subset\Z^2$.  This makes the ordering problem unusually explicit: a greedy
choice is a problem of selecting exponent vectors whose images spread through the cyclic subgroup
$\langle2,3\rangle$ with maximal $p$-adic separation.

The route proposed here has four precise stages.
\begin{enumerate}[leftmargin=2.2em]
\item \textbf{Compute the generalized factorials of $\mathcal U$.}  Express their $p$-adic
      valuations in terms of the orders of $2$ and $3$ modulo $p^m$ and the index of
      $\langle2,3\rangle$ in $(\Z/p^m\Z)^\times$.
\item \textbf{Interpolate $\Phi$ on a $p$-ordering.}  The Newton coefficients are divided
      differences of values $M^aA^b$.  Their required divisibility can be stated entirely in
      $\Z_p$.
\item \textbf{Convert coefficient divisibility to kernel inclusion.}  One seeks a theorem
      adapted to this sparse multiplicative set, showing that the special exponential form of
      the values forces the local condition in \cref{prop:kernelcriterion}.
\item \textbf{Use the density obstruction.}  If that condition holds at almost all primes,
      \cref{thm:localcommon} gives full local common exponents, while
      \cref{thm:densitygap} rules out both nonintegral rank branches.
\end{enumerate}

\begin{warning}
Bhargava's theory characterizes polynomial functions and integer-valued interpolation under
specific hypotheses; it does not state that an arbitrary monoid homomorphism is
congruence-preserving.  The proposed use is a coordinate system and divisibility framework, not a
citation that closes the problem.
\end{warning}

\begin{problem}[Sparse-monoid factorial problem]\label{prob:factorial}
Determine, for every odd prime $p$, the Bhargava factorial sequence of
$\mathcal U=\{2^a3^b:a,b\ge0\}$ in terms of the filtration
\[
       \ker\rho_p\supseteq\ker\rho_{p^2}\supseteq\cdots.
\]
Then characterize those pairs $(M,A)$ for which the map $2^a3^b\mapsto M^aA^b$ has Newton
coefficients satisfying the congruence-preserving divisibility bounds.
\end{problem}

This target is finite at every level, branch-sensitive through the Kummer relation, and directly
connected to the exact support criterion.  It is therefore better focused than a generic appeal
to integer-valued polynomial theory.

\section{Finite regulators, horizontal sieves, and the missing bridge}
\label{sec:regulator}

The exact density theorem solves the local distribution problem conditional only on multiplicative
rank.  What remains is to make the real identity communicate with those local characters.  This
section formulates that gap in three related ways.

\subsection{An archimedean tensor killed by one bilinear functional}

In the tensor square of the rational multiplicative group, consider
\[
       \tau=2\otimes A-3\otimes M
       \in\Gamma\otimes_{\Z}\Gamma.
       \tag{10.1}\label{eq:tau}
\]
The archimedean logarithm gives a bilinear functional
\[
       B_\infty(u\otimes v)=\log u\,\log v.
\]
The candidate equations say exactly
\[
       B_\infty(\tau)
       =\log2\,\log A-\log3\,\log M=0.
       \tag{10.2}\label{eq:archdet}
\]

At a prime $p\equiv1\pmod Q$, the residue character
$\lambda_{p,Q}:\Gamma\to R_Q$ gives the finite determinant
\[
 D_{p,Q}
 =\lambda_{p,Q}(2)\lambda_{p,Q}(A)
  -\lambda_{p,Q}(3)\lambda_{p,Q}(M).
 \tag{10.3}\label{eq:finitedet}
\]
This is the image of $\tau$ under the corresponding finite bilinear functional.  $Q$-coherence
implies $D_{p,Q}=0$.  At depth one, if the base vector
$(\lambda(2),\lambda(3))$ is nonzero, the converse also holds because two vectors over a field
are proportional exactly when their determinant vanishes.  Over $R_Q$ with $e>1$, determinant
vanishing is weaker because of zero divisors, which is why the cyclic-submodule count in
\cref{thm:rank4density,thm:rank3density} is the correct object.

The conjectural missing principle can now be stated without mentioning exponentials.

\begin{problem}[Finite-regulator principle]\label{prob:regulator}
Find an arithmetic reciprocity law, product formula, height inequality, or regulator construction
for tensors such as \eqref{eq:tau} with the following consequence: if a tensor built from positive
rational integers is killed by $B_\infty$ for the special reason
$\log_2 M=\log_3 A$, then its finite Kummer images
\eqref{eq:finitedet} vanish, or become coherent in the stronger cyclic-submodule sense, at a
set of primes whose density exceeds the exact rank-three and rank-four baselines.
\end{problem}

Ordinary product formulas apply to sums of valuations of one algebraic number; here the
archimedean expression is a product of logarithms and the finite data consists of power-residue
characters.  A solution may therefore require a genuinely higher object---for example a cup
product, a $K_2$ symbol, or an adelic regulator---rather than another linear form in logarithms.
The point of \cref{prob:regulator} is not to guess the correct cohomology in advance, but to give
it a numerical acceptance test: it must beat $P_3(q,e)$ and $P_4(q,e)$.

\subsection{Depth amplification}

Because both coherence probabilities tend to zero, even a weak lower bound uniform in depth
would settle the problem.

\begin{corollary}[Uniform-depth amplification criterion]\label{cor:depthamp}
Suppose there exist an admissible prime $q$, a constant $\delta>0$, and an unbounded set of
integers $e$ such that the candidate equations imply
\[
 \underline{\dens}\left(
   \{p\equiv1\pmod{q^e}:p\text{ is }q^e\text{-coherent}\}
   \;\middle|\;p\equiv1\pmod{q^e}
 \right)\ge\delta.
 \tag{10.4}\label{eq:depthamp}
\]
Then $x\in\Z$.
\end{corollary}
\begin{proof}
For a nonintegral candidate the exact conditional density is $P_3(q,e)$ or $P_4(q,e)$, both of
which tend to zero.  This contradicts \eqref{eq:depthamp} along the unbounded set of depths.
\end{proof}

This criterion is substantially weaker than congruence preservation.  It does not ask for local
common exponents at almost all primes, or even for density one.  Any positive amount of coherence
that survives arbitrarily deep $q$-power quotients is enough.

Possible sources of depth amplification include:
\begin{enumerate}[leftmargin=2.2em]
\item \textbf{A $q$-adic interpolation of the exponent.}  Construct, for many primes $p$, a
      compatible scalar in $\Z_q$ from the two equations and show its reductions solve
      \eqref{eq:coherent}.  The obstacle is that $x$ is a real number with no given $q$-adic
      incarnation.
\item \textbf{A reciprocity constraint on the tensor $\tau$.}  Show that
      $B_\infty(\tau)=0$ forces small finite regulator at all depths for a positive proportion
      of Frobenius elements.
\item \textbf{A congruence-preserving interpolation theorem.}  Use the special values of $\Phi$
      on a $p$-ordering of $\mathcal U$ to obtain local common exponents modulo $p^m$ for a
      set of primes with uniform lower density.
\item \textbf{An unlikely-intersection theorem in character space.}  The Kummer characters of
      $\Gamma$ form $R_Q^r$; coherence is a union of cyclic submodules.  One needs an
      arithmetic reason for the Frobenius orbit of the special candidate to hit this thin set
      more often than a generic rank-$r$ point.
\end{enumerate}

\subsection{Horizontal Frobenius sieve}

A full common exponent modulo $p$ implies $q$-coherence for every prime $q\mid p-1$.  Thus one
may impose many independent-looking Kummer filters horizontally, across distinct $q$.  For a
single admissible $q$, the filter rejects a proportion
\[
       \frac1{q-1}\bigl(1-P_r(q,1)\bigr)
\]
of all primes at the first $q$-level.  The sum of the leading rejection probabilities over $q$
diverges like $\sum_q1/q$.  This suggests that a large-sieve theorem for the compatible Kummer
representations should force the set
\[
 \mathcal P_{\mathrm{full}}
 =\{p:\exists k_p,
       M\equiv2^{k_p},\ A\equiv3^{k_p}\pmod p\}
 \tag{10.5}\label{eq:Pfull}
\]
to have density zero whenever $r>2$.

A zero-density theorem for \eqref{eq:Pfull} would strengthen the qualitative support theorem and
would be valuable independently of the conjecture.  It still would not prove Alaoglu--Erd\H{o}s,
because no current argument places almost all primes in $\mathcal P_{\mathrm{full}}$.  Its value
is strategic: it converts any future positive-density local-exponent consequence of the real
equations into an immediate contradiction.

The technical issue is entanglement.  Kummer fields for different $q$ share cyclotomic and
occasionally radical subfields, so one cannot multiply the one-prime probabilities naively.
The explicit entanglement-group theory of Perucca--Sgobba--Tronto is the natural tool for making
the horizontal sieve rigorous.

\begin{problem}[Horizontal common-exponent sieve]\label{prob:horizontalsieve}
Prove that if $\rank\langle2,3,M,A\rangle>2$, then the prime set
$\mathcal P_{\mathrm{full}}$ in \eqref{eq:Pfull} has natural density zero, with an effective
upper-bound sieve expressed through the Kummer entanglement invariants of $\Gamma$.
\end{problem}

\subsection{Relation recovery as a computational attack on rank three}

The rank-three branch is more structured and may be attacked separately.
For admissible $q$, exact residue logs at sufficiently many primes identify the hyperplane
$\bm c^\perp\subset\F_q^4$.  Repeating this for several $q$ reconstructs the primitive relation
$\bm c$, after which the real exponent has the fractional-linear expression
\[
       x=-\frac{c_1+c_2\theta}{c_3+c_4\theta},
       \qquad \theta=\log_2 3.
       \tag{10.6}\label{eq:mobiusx}
\]
The integrality conditions become the pair
\[
  2^{-\frac{c_1+c_2\theta}{c_3+c_4\theta}}\in\Z,
  \qquad
  3^{-\frac{c_1+c_2\theta}{c_3+c_4\theta}}\in\Z.
\]
This does not remove the four-exponentials wall, but it reduces the branch to a countable family
indexed by primitive integer matrices with nonzero determinant.  A branch-specific computation
can therefore prioritize relations of small height, test all local pattern consequences, and
search for additional congruences forced by the valuation supports of $M$ and $A$.

A useful formal target is an effective lower bound for the height of $\bm c$ in terms of $M,A$,
or a theorem excluding all relations with $|\Delta|$ in a prescribed range.  Such a theorem
would combine well with exact relation recovery modulo many $q$.

\section{Synthesis: what would now constitute a proof}
\label{sec:synthesis}

The new local theory organizes several sufficient routes by strength.

\begin{longtable}{p{0.25\textwidth}p{0.46\textwidth}p{0.20\textwidth}}
\toprule
Route & Required statement & Why it concludes\\
\midrule
\endfirsthead
\toprule
Route & Required statement & Why it concludes\\
\midrule
\endhead
Full congruence preservation &
$u-v\mid\Phi(u)-\Phi(v)$ for all $u,v\in\mathcal U$ &
\cref{prop:CP} and the support theorem\\[4pt]
One-dimensional divisibility &
$2^n-1\mid M^n-1$ for all sufficiently large $n$ &
Order support for $2,M$\\[4pt]
Almost-everywhere local exponent &
For almost all $p$, one $k_p$ satisfies
$M\equiv2^{k_p}$ and $A\equiv3^{k_p}\pmod p$ &
Implies $2\preccurlyeq M$ and $3\preccurlyeq A$\\[4pt]
Fixed-depth density excess &
For some admissible $(q,e)$, coherence density exceeds both exact baselines &
\cref{cor:quantcriterion}\\[4pt]
Uniform depth &
A positive lower coherence density survives unbounded $e$ &
\cref{cor:depthamp}\\[4pt]
Finite regulator &
$B_\infty(\tau)=0$ forces enough finite determinant or cyclic-submodule vanishing &
Reduces to either density criterion\\
\bottomrule
\end{longtable}

The additive results supply a complementary exclusion table.
\begin{longtable}{p{0.30\textwidth}p{0.31\textwidth}p{0.30\textwidth}}
\toprule
Input used & Maximum modulus obtained & Obstruction\\
\midrule
\endfirsthead
\toprule
Input used & Maximum modulus obtained & Obstruction\\
\midrule
\endhead
All primitive equations $u+v=w$ in $\mathcal U$ &
$C=\gcd(M-2,A-3)$ &
Fixed in $n$\\[4pt]
All $P(2,3)=0$ with $P\in\Z[X,Y]$ &
The same $C$ &
Universal ideal collapse\\[4pt]
All $P(2^n,3^n)=0$ at time $n$ &
$C_n=\gcd(M^n-2^n,A^n-3^n)$ &
$\log C_n=o(n)$\\[4pt]
Support target &
Needs to carry divisors of size about $e^{n\log2}$ &
Exponential gap\\
\bottomrule
\end{longtable}

The synthesis is therefore sharp.  The local side is not vague: its generic distributions and
thresholds are exact.  The additive side is not merely unsuccessful: a full ideal theorem and an
$S$-unit gcd estimate explain its ceiling.  The remaining problem is one specific transfer:
produce non-generic finite-place coherence from the special real logarithmic dependence.

\section{Lean formalization blueprint}
\label{sec:lean}

A useful formalization should separate elementary algebra, finite combinatorics, and deep imported
number theory.  The following module graph minimizes trust boundaries.

\subsection{Kernel-checkable elementary layer}

The following statements can be formalized in Lean with Mathlib and no analytic number theory.
\begin{enumerate}[leftmargin=2.2em]
\item \textbf{Candidate normalization.}  Rational candidate implies integral; rank-two exclusion
      assuming the transcendence of $\log_2 3$ as an external theorem.
\item \textbf{Relation algebra.}  Primitive relation lattice in rank three and
      $\Delta\ne0$ for a nonintegral candidate.
\item \textbf{Finite-ring content.}  Cardinality of vectors of exact $q$-adic content in
      $(\Z/q^e\Z)^2$ and size of the cyclic submodule they generate.
\item \textbf{Rank-four count.}  The identity
\[
  1+\sum_{t=1}^e(q^{2t}-q^{2t-2})q^t
  =1+\frac{q(q^2-1)(q^{3e}-1)}{q^3-1}.
\]
\item \textbf{Rank-three count.}  Invertible change of variables under $q\nmid\Delta$ and
\[
  1+\sum_{t=1}^e(q^{2t}-q^{2t-1})
  =1+\frac{q(q^{2e}-1)}{q+1}.
\]
\item \textbf{Pattern law.}  The recurrence and closed formula for $N_s(q)$, including the
      forbidden one-nonzero patterns.
\item \textbf{Additive skeleton.}  The elementary factorization proof of the four primitive
      triples in \eqref{eq:fourtriples}.
\item \textbf{Defect identities and ideals.}  Equations \eqref{eq:d13}--\eqref{eq:d18}, the
      kernel $(X-2,Y-3)$, and the universal polynomial collapse.
\item \textbf{Dynamic divisibility.}  $C_m\mid C_n$ when $m\mid n$.
\item \textbf{Local cyclic groups.}  Kernel inclusion and the common-exponent criterion for
      cyclic unit groups modulo odd prime powers.
\end{enumerate}

\subsection{External theorem interfaces}

The deep inputs should appear as explicit assumptions or imported theorems with narrow signatures.
\begin{description}[style=nextline,leftmargin=2em]
\item[Kummer realization]
For admissible $q$, identify the Galois group with the quotient of the exponent lattice computed
in \cref{lem:kummermodule}.
\item[Chebotarev]
For a finite abelian extension, assign density $|E|/|G|$ to a union of Frobenius elements, then
prove the scalar-invariant descent from prime ideals of $K_Q$ to rational primes.
\item[Support theorem]
Use the Corrales--Rodrig\'a\~nez--Schoof implication in the one-dimensional torus.
\item[$S$-unit gcd]
Package only the asymptotic statement needed for \cref{thm:Cnsubexp}; handle the finite support
primes by formalized LTE lemmas.
\end{description}

This separation would already allow a machine-checked conditional theorem:
\emph{assuming the four external interfaces, the exact density gap and additive no-go results
follow}.  Full formalization of Chebotarev or the Subspace Theorem is not a prerequisite for
checking the novel finite combinatorics.

\subsection{Computational certificates}

For fixed $(q,e)$, the theoretical numerator identities can be certified by exhaustive
enumeration in \texttt{ZMod}$(q^e)$ for small parameters.  Prime experiments should be represented as
independent certificates containing:
\begin{itemize}[leftmargin=2em,itemsep=1pt]
\item the prime list or a deterministic primality certificate;
\item a primitive $Q$th root in each finite field;
\item the four residue logs;
\item the coherence witness $k$ when it exists;
\item the relation check in rank three.
\end{itemize}
The proof does not depend on these certificates, but they are excellent regression tests for a
formal development.

\section{Prioritized research program}
\label{sec:program}

The following order maximizes the chance of converting the new framework into a proof rather than
another collection of analogies.

\begin{enumerate}[leftmargin=2.2em]
\item \textbf{Formalize the two coherence counts.}  They are short, exact, and likely to reveal
      any hidden issue with the prime-power content stratification.
\item \textbf{Prove the horizontal zero-density theorem} in \cref{prob:horizontalsieve}.  This
      is a natural standalone theorem about reductions of two points on $\mathbb G_m^2$ and can
      use existing entanglement machinery.
\item \textbf{Compute $p$-orderings of the $3$-smooth monoid.}  Begin with primes for which
      $\langle2,3\rangle$ has small index modulo $p^m$; search for a stable generalized-factorial
      formula.
\item \textbf{Test candidate-value divided differences symbolically.}  Express the first several
      Newton coefficients as Laurent polynomials in $M,A$ and factor their fixed divisors.  The
      additive collapse theorem predicts which factors are spurious and which could encode orders.
\item \textbf{Develop a finite regulator.}  Search first for a depth-one reciprocity law for
      $D_{p,q}$, where the determinant criterion is linear algebra over a field.  Only after a
      genuine density excess appears should the construction be lifted to $q^e$.
\item \textbf{Exploit the rank-three pattern detector.}  Combine relation reconstruction with
      valuation support to exclude low-height primitive relations; automate exact searches with
      proof-producing certificates.
\item \textbf{Maintain an adversarial claim ledger.}  Any proposed bridge must explicitly state
      whether it yields infinitely many coherent primes, positive density, density one, or uniform
      depth.  Only the last three can interact decisively with the exact baselines.
\end{enumerate}

\section{Conclusion}

This investigation does not close the Alaoglu--Erd\H{o}s conjecture, but it changes the shape of
the remaining work.

First, the conjecture has an exact support-theoretic endpoint: it is enough to force
$2\preccurlyeq M$ or $3\preccurlyeq A$.  Congruence preservation of the candidate monoid map is
a concrete sufficient mechanism, and locally it is equivalent to the existence of one common
finite exponent for the pairs $(2,3)$ and $(M,A)$.

Second, the finite-place behavior of every hypothetical nonintegral candidate is now quantified.
At prime-power depth $q^e$, common-exponent coherence has one of two exact densities, determined
solely by whether the multiplicative rank is three or four.  Both decay like $q^{-e}$, while an
integral exponent has density one.  The rank-three relation also leaves a visible fingerprint of
forbidden $q$th-power patterns.

Third, the obvious additive route has been exhausted in a rigorous sense.  The entire additive
skeleton of the $3$-smooth monoid consists of four primitive triples; all their defects, and in
fact all integral polynomial defects at $(2,3)$, generate only
$\gcd(M-2,A-3)$.  Dynamic versions are subexponential by the $S$-unit gcd theorem.  Additional
polynomial identities cannot bridge an exponential support gap unless they retain new
prime-dependent order data.

The most promising overlooked direction is therefore neither another interpolation determinant
nor another generic transcendence estimate.  It is a local--global regulator or sparse-monoid
interpolation theorem that converts the exact real logarithmic determinant into a non-generic
supply of finite common exponents.  The acceptance threshold is explicit: beat
$P_3(q,e)$ and $P_4(q,e)$ at one depth, or retain any positive density through unbounded depth.
That is a sharply falsifiable research objective and a tractable target for formal proof
engineering.

\appendix

\section{Derivation checks for the coherence numerators}

For reference, the rank-four numerator may be rearranged as
\[
\begin{aligned}
 N_4(q,e)
 &=1+\sum_{t=1}^e(q^{2t}-q^{2t-2})q^t\\
 &=1+(q^2-1)\sum_{t=1}^e q^{3t-2}\\
 &=1+q(q^2-1)\frac{q^{3e}-1}{q^3-1}.
\end{aligned}
\]
Its leading contribution comes from primitive base vectors ($s=0$, $t=e$):
\[
       (q^{2e}-q^{2e-2})q^e
       =(1-q^{-2})q^{3e}.
\]
After division by $q^{4e}$ this gives
\[
  (1-q^{-2})q^{-e}
  \frac{1}{1-q^{-3}}+O(q^{-4e})
  =\frac{q(q+1)}{q^2+q+1}q^{-e}+O(q^{-4e}).
\]

Similarly,
\[
\begin{aligned}
 N_3(q,e)
 &=1+\sum_{t=1}^e(q^{2t}-q^{2t-1})\\
 &=1+q(q-1)\frac{q^{2e}-1}{q^2-1}\\
 &=1+\frac{q(q^{2e}-1)}{q+1},
\end{aligned}
\]
so
\[
       P_3(q,e)=\frac{q}{q+1}q^{-e}+O(q^{-3e}).
\]

\section{A compact reproducibility script}
\label{app:code}

The following script reproduces the pairwise, status-pattern, and coherence experiments.  It uses
only Python's exact integers.  The depth-$25$ loop is intentionally separated because it uses a
larger prime bound.

\begin{lstlisting}[style=pythonstyle]
from __future__ import annotations
from collections import Counter
from math import isqrt


def primes_upto(n: int) -> list[int]:
    sieve = bytearray(b"\x01") * (n + 1)
    sieve[:2] = b"\x00\x00"
    for p in range(2, isqrt(n) + 1):
        if sieve[p]:
            sieve[p*p:n+1:p] = b"\x00" * (((n - p*p)//p) + 1)
    return [i for i, flag in enumerate(sieve) if flag]


def residue_logs(vals: tuple[int, ...], p: int, Q: int, q: int):
    n = (p - 1)//Q
    omega = None
    for g in range(2, 200):
        candidate = pow(g, n, p)
        if pow(candidate, Q//q, p) != 1:  # exact order Q
            omega = candidate
            break
    if omega is None:
        raise RuntimeError(("no root", p, Q))
    table = {}
    cur = 1
    for j in range(Q):
        table[cur] = j
        cur = cur * omega % p
    return tuple(table[pow(v % p, n, p)] for v in vals)


def coherent(t: tuple[int, int, int, int], Q: int) -> bool:
    u, v, m, a = t
    return any((k*u - m) % Q == 0 and (k*v - a) % Q == 0
               for k in range(Q))


def predicted(q: int, e: int, rank: int) -> float:
    Q = q**e
    if rank == 4:
        numerator = 1 + q*(q*q - 1)*(q**(3*e) - 1)//(q**3 - 1)
        return numerator/Q**4
    numerator = 1 + q*(q**(2*e) - 1)//(q + 1)
    return numerator/Q**3


def coherence_run(limit: int, q: int, e: int):
    Q = q**e
    primes = primes_upto(limit)
    examples = [
        ("rank4", (2, 3, 7, 11), 4),
        ("rank3", (2, 3, 7, 42), 3),
    ]
    for name, vals, rank in examples:
        total = good = relation_bad = 0
        for p in primes:
            if p % Q != 1 or any(v % p == 0 for v in vals):
                continue
            t = residue_logs(vals, p, Q, q)
            total += 1
            good += coherent(t, Q)
            if rank == 3 and (t[0] + t[1] + t[2] - t[3]) % Q:
                relation_bad += 1
        print(name, "Q", Q, "total", total, "good", good,
              "frequency", good/total,
              "predicted", predicted(q, e, rank),
              "relation_bad", relation_bad)


def is_qth_power(a: int, p: int, q: int) -> bool:
    return pow(a % p, (p - 1)//q, p) == 1


def status_run(limit: int = 5_000_000, q: int = 5):
    primes = primes_upto(limit)

    pair = Counter()
    for p in primes:
        if p % q == 1 and p not in (2, 7):
            pair[(is_qth_power(2, p, q), is_qth_power(14, p, q))] += 1
    print("pair", pair)

    for name, vals in [("rank4", (2, 3, 7, 11)),
                       ("rank3", (2, 3, 7, 42))]:
        counts = Counter()
        for p in primes:
            if p % q != 1 or any(v % p == 0 for v in vals):
                continue
            pattern = tuple(int(is_qth_power(v, p, q)) for v in vals)
            counts[pattern] += 1
        print(name, counts)


if __name__ == "__main__":
    status_run()
    coherence_run(5_000_000, 5, 1)
    coherence_run(10_000_000, 5, 2)
\end{lstlisting}

\section{Source and novelty audit}

The nonduplication audit used the following checks on the supplied source.
\begin{itemize}[leftmargin=2em]
\item Section headings and labels were extracted programmatically; the source contains a unified
      report, two large continuations, twenty-seven shorter reports, and multiple appendices.
\item Exact searches were run for the terminology central to this continuation.  No hits were
      found for the support theorem authors, Chebotarev, the support problem, discrete logarithms,
      or congruence preservation.  A normalized $20$-token shingle comparison found no substantive
      common passage; its sole hit was shared theorem-environment boilerplate in the preambles.
\item Generic Kummer theory, valuation lattices, and the Bugeaud--Corvaja--Zannier theorem do
      occur in the source.  The present use differs: Kummer theory supplies exact residue-vector
      densities; the gcd theorem proves the dynamic additive-anchor ceiling.
\item The additive classification, defect ideal, prime-power coherence counts, relation-pattern
      law, and local cyclic common-exponent criterion were derived independently for this report.
\end{itemize}

\section{Bibliography and public artifacts}

\begin{thebibliography}{99}

\bibitem{AE1944}
L.~Alaoglu and P.~Erd\H{o}s,
\emph{On highly composite and similar numbers},
Transactions of the American Mathematical Society \textbf{56} (1944), 448--469.
\href{https://doi.org/10.1090/S0002-9947-1944-0011087-2}{doi:10.1090/S0002-9947-1944-0011087-2}.

\bibitem{CRS1997}
C.~Corrales-Rodrig\'a\~nez and R.~Schoof,
\emph{The support problem and its elliptic analogue},
Journal of Number Theory \textbf{64} (1997), 276--290.
\href{https://doi.org/10.1006/jnth.1997.2114}{doi:10.1006/jnth.1997.2114}.

\bibitem{Bhargava1997}
M.~Bhargava,
\emph{$P$-orderings and polynomial functions on arbitrary subsets of Dedekind rings},
Journal f\"ur die reine und angewandte Mathematik \textbf{490} (1997), 101--128.
\href{https://doi.org/10.1515/crll.1997.490.101}{doi:10.1515/crll.1997.490.101}.

\bibitem{BCZ2003}
Y.~Bugeaud, P.~Corvaja, and U.~Zannier,
\emph{An upper bound for the G.C.D. of $a^n-1$ and $b^n-1$},
Mathematische Zeitschrift \textbf{243} (2003), 79--84.
\href{https://doi.org/10.1007/s00209-002-0449-z}{doi:10.1007/s00209-002-0449-z}.

\bibitem{CZ2004}
P.~Corvaja and U.~Zannier,
\emph{A lower bound for the height of a rational function at $S$-unit points},
Monatshefte f\"ur Mathematik \textbf{144} (2005), 203--224; preprint 2003.
\href{https://doi.org/10.1007/s00605-004-0273-0}{doi:10.1007/s00605-004-0273-0}.

\bibitem{PST2020}
A.~Perucca, P.~Sgobba, and S.~Tronto,
\emph{Explicit Kummer theory for the rational numbers},
International Journal of Number Theory \textbf{16} (2020), no.~10, 2213--2231.
\href{https://doi.org/10.1142/S1793042120501146}{doi:10.1142/S1793042120501146}.

\bibitem{PST2022}
A.~Perucca, P.~Sgobba, and S.~Tronto,
\emph{Kummer theory for number fields via entanglement groups},
Manuscripta Mathematica \textbf{169} (2022), 251--270.
\href{https://doi.org/10.1007/s00229-021-01328-0}{doi:10.1007/s00229-021-01328-0}.

\bibitem{Perucca2011}
A.~Perucca,
\emph{On the reduction of points on abelian varieties and tori},
International Mathematics Research Notices \textbf{2011} (2011), 293--308.
\href{https://doi.org/10.1093/imrn/rnq080}{doi:10.1093/imrn/rnq080}.

\bibitem{DebryPerucca2013}
C.~Debry and A.~Perucca,
\emph{Reductions of subgroups of the multiplicative group},
arXiv:1312.6620 (2013).

\bibitem{BruinPerucca2018}
P.~Bruin and A.~Perucca,
\emph{Reductions of points on algebraic groups, II},
arXiv:1802.08527 (2018).

\bibitem{GelfondSchneider}
A.~O.~Gelfond,
\emph{Transcendental and Algebraic Numbers}, Dover, 1960;
and Th.~Schneider, \emph{Einf\"uhrung in die transzendenten Zahlen}, Springer, 1957.

\bibitem{Chebotarev}
J.~Neukirch,
\emph{Algebraic Number Theory}, Springer, 1999, Chapter VII.

\bibitem{JacobianLean2026}
P.~Nogueira Grossi,
\emph{An Independent Lean 4 Verification of the Alp\"oge--Fable Counterexample},
Zenodo, version 1.0, 23 July 2026.
\url{https://zenodo.org/records/21514514}.

\bibitem{Gao2026}
S.~Gao,
\emph{Counterexamples to the Jacobian conjecture in dimensions greater than two},
arXiv:2608.00222 (2026).

\bibitem{ICARM273}
NSF Institute for Computer-Aided Reasoning in Mathematics,
\emph{Elliptic Curve Rank Leaderboard, curve 273}, accessed 22 August 2026.
\url{https://elliptic-rank.icarm.cloud/curve/273}.

\bibitem{OEISHadamard}
OEIS Foundation,
\emph{A007299 and A019442: Hadamard-matrix existence records},
updates of 12--13 August 2026.
\url{https://oeis.org/A007299}; \url{https://oeis.org/A019442}.

\end{thebibliography}

\end{document}
